\documentclass[12pt,letterpaper]{article}
% Plantilla: la del Syllabus de Cálculo Integral. Compilar con pdflatex.
\usepackage[spanish,es-nodecimaldot]{babel}
\usepackage[utf8]{inputenc}
\usepackage[T1]{fontenc}
\usepackage{lmodern}
\usepackage{geometry}
\usepackage{array}
\usepackage{booktabs}
\usepackage{longtable}
\usepackage{tabularx}
\usepackage{enumitem}
\usepackage{hyperref}
\usepackage{xcolor}
\usepackage{titlesec}
\usepackage{setspace}
\usepackage{amsmath, amssymb}
\geometry{letterpaper,left=2.5cm,right=2.5cm,top=2.5cm,bottom=2.5cm}
\setstretch{1.12}
\setlength{\parskip}{0.6em}
\setlength{\parindent}{0pt}
\hypersetup{colorlinks=true,linkcolor=black,urlcolor=blue,citecolor=black}
\titleformat{\section}{\normalfont\Large\bfseries}{\thesection.}{0.5em}{}
\titleformat{\subsection}{\normalfont\large\bfseries}{\thesubsection.}{0.5em}{}
\renewcommand{\arraystretch}{1.25}
\begin{document}
\begin{center}
    {\Large \textbf{Universidad Nacional de Colombia}}\\
    {\large \textbf{Sede Manizales}}\\
    {\large \textbf{Departamento de Matemáticas y Estadística}}\\[1em]
    {\LARGE \textbf{Syllabus del curso}}\\[0.4em]
    {\LARGE \textbf{Taller en Ciencia de Datos y Aprendizaje Automático}}\\[0.4em]
    {\large veinte proyectos con datos reales}\\[1em]
\end{center}

\section{Encabezado institucional}
\begin{tabularx}{\textwidth}{>{\bfseries}p{4.5cm}X}
Institución: & Universidad Nacional de Colombia, Sede Manizales.\\
Facultad: & Ciencias Exactas y Naturales.\\
Departamento: & Matemáticas y Estadística.\\
Curso: & Taller en Ciencia de Datos y Aprendizaje Automático.\\
Programa: & 4345 Especialización en Estadística.\\
Tipo de asignatura: & Elegible.\\
Semestre: & 2026-II.\\
Créditos: & 4.\\
Sesiones: & Doce sesiones de cuatro horas, del 20 de noviembre al 14 de diciembre de 2026.\\
Modalidad: & Teórico-práctica.\\
Asistencia mínima: & 80\,\%.\\
Herramientas: & Python en Google Colab, GitHub y Claude como asistente de IA del curso.\\
Profesor: & Carlos Manuel Orrego Franco.\\
Prerrequisitos: & Conocimientos de aprendizaje automático y ciencia de datos.\\
Repositorio: & \url{https://github.com/cmorregof/taller-cd-ml-2026-2}\\
\end{tabularx}

\section{Presentación del curso}
Este taller cierra la Especialización. Su propósito es que cada estudiante salga capaz de llevar un proyecto de datos de principio a fin: formular la pregunta, conseguir y limpiar el dato, modelar, y defender una decisión ante alguien que no es técnico.

El programa oficial propone un solo proyecto por equipo durante el semestre. Este taller conserva sus seis resultados de aprendizaje y cambia la forma de practicarlos. En lugar de recorrer el ciclo de vida una vez, lo recorremos veinte veces en proyectos cortos, porque un oficio se aprende repitiéndolo.

Cada clase tiene dos proyectos. El primero lo resuelve el profesor en vivo, en un cuaderno de Colab. El segundo es su gemelo: la misma técnica en otro mundo de datos, y lo resuelve cada estudiante por su cuenta. Los datos son reales en todos los casos y se descargan de su fuente oficial. Hay precios de alimentos, mercado eléctrico, salud pública, imágenes médicas y satelitales, texto, audio, deporte y contratación pública.

No hace falta ser programador. El curso entrena tres habilidades: leer código, completarlo, y saber qué pedirle a una IA y cómo auditar lo que responde.

\section{Objetivo general}
Diseñar, ejecutar y comunicar un proyecto aplicado de Ciencia de Datos o de Machine Learning, integrando todas las fases del ciclo de vida de un proyecto de análisis y modelado con datos.

\section{Resultados de aprendizaje}
Son los del programa oficial del curso. Al finalizar, el estudiante estará en capacidad de:
\begin{enumerate}[label=\arabic*.]
    \item Formular un problema analítico del mundo real, identificando objetivos claros, métricas de éxito y criterios de evaluación, de forma colaborativa.
    \item Seleccionar y preparar conjuntos de datos apropiados, aplicando técnicas de limpieza, transformación y exploración que permitan extraer valor y detectar patrones.
    \item Diseñar y ejecutar un flujo de trabajo reproducible para el desarrollo de un proyecto de Ciencia de Datos o Machine Learning, siguiendo principios de trazabilidad, documentación y buenas prácticas de codificación.
    \item Implementar y evaluar un modelo de análisis o predicción, utilizando métricas cuantitativas para validar su desempeño y tomar decisiones basadas en evidencia.
    \item Documentar y comunicar los resultados técnicos y estratégicos del proyecto, utilizando herramientas visuales y discursivas adecuadas a diferentes audiencias (técnicas y no técnicas).
    \item Reflexionar críticamente sobre los aspectos éticos, técnicos y operativos del proyecto, identificando oportunidades de mejora, riesgos asociados y posibles rutas de mantenimiento o escalamiento futuro.
\end{enumerate}

\section{Cómo funciona una clase}
Cada sesión dura cuatro horas. En la primera mitad el profesor resuelve un proyecto completo y los estudiantes lo corren y lo leen línea a línea. En la segunda mitad cada estudiante resuelve el gemelo. Al cerrar, el cuaderno queda en el portafolio de GitHub de cada uno.

Todos los cuadernos abren con una ficha que dice el objetivo pedagógico, los resultados de aprendizaje que trabaja, el uso permitido de la IA y el entregable. Después vienen cuatro secciones, que son los cuatro módulos del programa oficial y aparecen siempre en el mismo orden:
\begin{enumerate}[label=\arabic*.]
    \item \textbf{Formulación.} El problema, la métrica de éxito y una predicción escrita antes de ver el dato.
    \item \textbf{Datos.} La descarga desde la fuente oficial, la limpieza con su registro y la exploración.
    \item \textbf{Modelado.} Una línea base, varios modelos comparados y una evaluación que respeta el tiempo o la estructura del dato.
    \item \textbf{Comunicación.} La decisión en tres frases, la bitácora de decisiones y una reflexión sobre riesgos y mantenimiento.
\end{enumerate}

Cada mundo de datos trae tres preguntas. \emph{Predecir} pregunta qué va a pasar. \emph{Explicar} pregunta por qué pasa y qué no puede afirmar el modelo. \emph{Decidir} pregunta qué se hace con el resultado y cuánto cuesta equivocarse.

La ayuda se retira por etapas. En las clases 1 a 4 el cuaderno gemelo trae el código casi completo, con huecos por llenar. En las clases 5 a 8 trae solo los títulos de las secciones. En las clases 9 y 10 se parte de un repositorio casi vacío, como en un proyecto real.

\newpage
\section{Calendario}
\begin{longtable}{p{0.5cm}p{1.6cm}p{4.1cm}p{4.1cm}p{3.9cm}}
\toprule
\textbf{\#} & \textbf{Fecha} & \textbf{Resuelve el profesor} & \textbf{Resuelven los estudiantes} & \textbf{Frente}\\
\midrule
\endfirsthead
\toprule
\textbf{\#} & \textbf{Fecha} & \textbf{Resuelve el profesor} & \textbf{Resuelven los estudiantes} & \textbf{Frente}\\
\midrule
\endhead
1 & Vie 20 nov & El sancocho (SIPSA, DANE) & El precio de la luz (XM) & Formular, línea base y validación temporal.\\
2 & Sáb 21 nov & El ICFES (Saber 11) & Dengue (SIVIGILA, Bucaramanga y Tuluá) & Limpieza, sesgos y métrica.\\
3 & Mar 24 nov & La factura de Londres (London Datastore) & El cielo de Caldas (IDEAM y UNGRD) & Datos grandes y versionamiento de datos.\\
4 & Vie 27 nov & El Mundial de Qatar (StatsBomb) & Las calles de Bogotá (siniestros viales) & Desbalance y calibración. Simulacro de la prueba técnica.\\
5 & Sáb 28 nov & \textbf{Prueba técnica} en la primera hora. Lunares (HAM10000) & La roya (RoCoLe) & Imágenes, transferencia y fuga de datos. Se asigna el caso ejecutivo.\\
6 & Mar 1 dic & \textbf{Caso ejecutivo.} La tierra desde arriba (EuroSAT y Caldas) & Fauna de la Orinoquía, en trabajo autónomo & Comunicación, revisión por pares y cambio de dominio.\\
7 & Vie 4 dic & Lo que dicen los turistas (Rest-Mex) & La plata pública (SECOP II) & Texto.\\
8 & Sáb 5 dic & El electrocardiograma (PTB-XL) & Los sonidos del Magdalena (BirdCLEF) & Señales y audio.\\
9 & Mié 9 dic & El noticiero de un minuto & El parte de incendios (NASA FIRMS) & API, despliegue y agentes.\\
10 & Vie 11 dic & La cifra que falta (JEP, CEV y HRDAG) & El mapa de vuelos (Aerocivil) & Incertidumbre, tablero y contar el resultado.\\
11 & Sáb 12 dic & \multicolumn{2}{p{8.2cm}}{Armado del proyecto final, cuatro horas de taller.} & Integración.\\
12 & 14 dic & \multicolumn{2}{p{8.2cm}}{\textbf{Sustentación final}, individual.} & Comunicación.\\
\bottomrule
\end{longtable}

\section{Política de uso de inteligencia artificial}
El asistente de IA del curso es Claude y su uso está permitido en clase. La regla es que la IA trabaja para el estudiante y el estudiante responde por todo lo que quede en su repositorio.

Esa regla se concreta en tres hábitos. El primero es \emph{leer}: entender cada línea antes de correrla. El segundo es \emph{pedir}: saber qué se le pide a la IA y cómo se le pide, porque una pregunta sin contexto recibe una respuesta genérica. El tercero es \emph{auditar}: revisar lo que la IA entrega antes de confiar en ello, en especial la forma en que evalúa un modelo.

La ficha de cada cuaderno dice qué se le puede pedir a la IA en ese proyecto y qué debe poder explicar el estudiante sin ella. En la prueba técnica, en el caso ejecutivo y en la sustentación final cada estudiante explica su trabajo sin ayuda.

\section{Esquema de evaluación}
\begin{table}[h!]
\centering
\begin{tabularx}{\textwidth}{Xr}
\toprule
\textbf{Componente} & \textbf{Ponderación}\\
\midrule
Portafolio: los proyectos de clase, resueltos y publicados en GitHub & 20\,\%\\
Prueba técnica escrita (sábado 28 de noviembre) & 20\,\%\\
Caso ejecutivo (martes 1 de diciembre) & 20\,\%\\
Sustentación final (14 de diciembre) & 40\,\%\\
\midrule
\textbf{Total} & \textbf{100\,\%}\\
\bottomrule
\end{tabularx}
\end{table}

\textbf{Portafolio.} Es el repositorio de GitHub de cada estudiante con los cuadernos gemelos. Un cuaderno cuenta cuando corre de principio a fin desde un clon limpio y tiene llenas la predicción escrita, la decisión en tres frases y la bitácora. Un repositorio que no corre no se califica.

\textbf{Prueba técnica.} Dura 75 minutos, es individual, escrita a mano y sin dispositivos. Simula la segunda etapa de un proceso de selección para un cargo de analista de datos. Cada estudiante recibe una carpeta con un análisis terminado, hecho con ayuda de un asistente de IA, y debe encontrar sus errores, proponer la siguiente comprobación, reescribir la pregunta que se le hizo a la IA y responderle en tres frases a quien decide. El viernes 27 hay un simulacro corto que no tiene nota.

\textbf{Caso ejecutivo.} El sábado 28 cada estudiante recibe un mundo de datos propio, con sus tres preguntas y un material inicial. El martes 1 de diciembre lo expone en diez minutos ante el grupo, como se le presentaría a una gerencia. Cada estudiante evalúa además a dos compañeros con una rúbrica corta.

\textbf{Sustentación final.} Cada estudiante formula, construye y defiende un proyecto propio. La sesión del sábado 12 de diciembre se dedica completa a armarlo.

\section{Recursos y materiales}
Todo el material está en el repositorio del curso. Cada cuaderno se abre en Google Colab con un clic desde la tabla del archivo \texttt{README}, así que no hace falta instalar nada. Los datos no están en el repositorio: cada cuaderno los descarga de su fuente oficial y deja anotado de dónde salieron y cuándo.

Para la primera sesión cada estudiante necesita un portátil, una cuenta de Google y una cuenta de GitHub.

\section{Bibliografía}
Es la del programa oficial del curso.

Jurney, R. (2017). \textit{Agile Data Science 2.0}. O'Reilly Media.

Géron, A. (2019). \textit{Hands-On Machine Learning with Scikit-Learn, Keras, and TensorFlow}. O'Reilly.

Cordeiro, R., Alves, I., Alves, S., \& Goldman, A. (2024). Being Agile in a Data Science Project. En P. Kruchten \& P. Gregory (eds.), \textit{Agile Processes in Software Engineering and Extreme Programming: Workshops}. Lecture Notes in Business Information Processing, vol. 489. Springer. \url{https://doi.org/10.1007/978-3-031-48550-3_6}

Provost, F., \& Fawcett, T. (2013). \textit{Data Science for Business}. O'Reilly.

Kuhn, M., \& Johnson, K. (2013). \textit{Applied Predictive Modeling}. Springer.

Data Science PM. (2024). \textit{Data Science Methodologies and Frameworks Guide}. \url{https://www.datascience-pm.com/data-science-methodologies/}
\end{document}
